\section{Proposed FPGA Decoder Accelerator}
\label{sec:accelerator}

\subsection{Architecture Overview}
Figure~\ref{fig:accelerator} shows the principal architecture organization.
The accelerator consists of a graph sequencer, a weighted-layer controller,
an INT8 systolic array, three auxiliary engines, and a shared memory access
path. The sequencer issues an operation descriptor and advances after the
selected engine completes. A complete decoder job therefore uses the same
hardware across its layers, rather than instantiating one engine per graph
node. The current cluster activates one operator engine at a time.

Weights, inputs, intermediate features, and outputs are backed by HBM. Local
M20K memories and registers supply operands to the compute array. This
distinguishes the implementation from a fully on-chip, layer-unrolled
streaming network: local computation is pipelined, but intermediate graph
edges may require HBM reads and writes. The design description and bandwidth
evaluation retain these transfers explicitly.

\subsection{Shared INT8 Systolic Compute Array}
The current RTL defines a $104\times32$ array: rows correspond to spatial
positions and columns to output channels. Each physical processing element
(PE) consumes two activation--weight pairs along the reduction dimension and
maintains an INT32 partial sum. A reduction tile contains up to 32 elements;
successive tiles accumulate into the same output tile. Valid-row,
valid-column, and valid-reduction counts handle boundary tiles.

For paired operands, the conceptual PE update is
\begin{equation}
  p_{r,c}\leftarrow p_{r,c}
       +a_{r,2k}w_{2k,c}+a_{r,2k+1}w_{2k+1,c}.
  \label{eq:pe}
\end{equation}
The implementation uses the DSP sum-of-two-products mode, with signed
operands and wider accumulation~\cite{intels10mxdsp}. Its structural maximum
is $104\times32\times2=6{,}656$ INT8 products/cycle. This count excludes
loading, pipeline fill, partial tiles, stalls, and result draining; sustained
throughput must be measured separately.

Valid and completion tokens follow operands through neighboring PEs.
After computation, results are drained column by column in spatial
order. The RTL uses local token propagation and a scan-style drain to avoid
broadcasting an array-wide result address to every accumulator. The timing
and resource impact of this organization remains part of the implementation
evaluation.

\subsection{Layer-Weight and Activation Buffering}
The weighted-layer controller fetches a layer's INT8 weights into a local
M20K cache and reuses them across spatial tiles. Two register tiles stage
weights for the array. Two activation-strip buffers support fetching the next
width-contiguous region while the current region is processed. Boundary
padding and address generation are performed outside the PE array.

The controller packs results into 64-bit memory words after bias and
requantization. Burst requests and ready/valid responses must be respected
throughout the memory path. The evaluation will report effective bandwidth
and stall cycles at this interface, rather than infer performance from HBM
capacity or the array's arithmetic peak.
